\subsection{Causal Basis Block}
To estimate feature expectations for the front-door adjustment, we employ a \textbf{subspace projection}~\cite{strang2022introduction} mechanism,
where input features are projected onto a learnable low-dimensional subspace spanned by a set of basis vectors.
\subsubsection{Expectations Estimation} 
\label{sec:coefficients_estimation}
Given an input feature representation $\mathcal{X} \in \mathbb{R}^{B \times N \times C}$, 
we introduce a set of learnable bases
$\mathcal{B} = [b_1, \dots, b_K] \in \mathbb{R}^{K \times C}$ ($K < C$),
which span a low-rank subspace.

For each feature vector, we seek coefficients
$\mathcal{C} \in \mathbb{R}^{B \times N \times K}$
that minimize the reconstruction error in this subspace:
\begin{equation}
\arg\min \sum_{i=1}^{N} \| \mathcal{X}_i - (\mathcal{C} \mathcal{B})_i \|^2,
\end{equation}
where the sum is taken over the $N$ samples in the batch.
This objective ensures that the projected features $\mathcal{C} \mathcal{B}$
are the closest approximation to $\mathcal{X}$ within the low-rank subspace spanned by $\mathcal{B}$.

The above minimization admits a closed-form least-squares solution:
\begin{equation}
\mathcal{C} = (\mathcal{B}^\top \mathcal{B})^{-1} \mathcal{B}^\top \mathcal{X}.
\end{equation}
Here, $(\mathcal{B}^\top \mathcal{B})^{-1}$ accounts for the fact that the bases may not remain orthogonal during training.
The reconstructed features
\begin{equation}
\mathcal{C} \mathcal{B}
= (\mathcal{B} (\mathcal{B}^\top \mathcal{B})^{-1} \mathcal{B}^\top)\mathcal{X}
\end{equation}
thus correspond to the projection of $\mathcal{X}$ onto the low-rank subspace.

Crucially, minimizing the squared reconstruction error over all feature vectors
can be interpreted as computing the expectation in the squared-error sense:
\begin{equation}
\mathbb{E}[X] 
= \arg\min \; \mathbb{E}\big[\| X - Y \|^2\big]
\approx \arg\min \sum_{i=1}^{N} \| X_i - Y_i \|^2,
\end{equation}
where the sum over $N$ samples provides an empirical estimate of the expectation.

Accordingly, we define the subspace projection function, which provides a low-rank approximation of the feature expectation:

\begin{equation}
f_{\mathrm{sp}}(\mathcal{X})
\triangleq \mathcal{C} \mathcal{B}
\approx \mathbb{E}[\mathcal{X}]
\in \mathbb{R}^{B \times N \times C}.
\label{equ:subspace_projection}
\end{equation}


Note that the entire procedure is \textbf{differentiable} and can be optimized end-to-end with the network.
At inference, the term $\mathcal{B} (\mathcal{B}^{\top} \mathcal{B})^{-1} \mathcal{B}^{\top}$ can be precomputed and absorbed into a single $C \times C$ matrix, reducing computation and simplifying deployment.



To further capture the generalized information of the input features,
we introduce a set of learnable \textit{Sample Queries}
$\mathcal{Q}_s \in \mathbb{R}^{1 \times S \times C}$,
which play a role analogous to the object queries in DETR~\cite{carion2020end}
and Mask2Former~\cite{cheng2022masked}.
During training, these queries implicitly aggregate global representations
across samples, thereby facilitating a more stable estimation of feature expectations.

The \textit{Sample Queries} interact with the input features to compute a set of correlation maps,
which are aggregated into an expected spatial weighting map:
\begin{equation}
\resizebox{0.433\textwidth}{!}{%
$\displaystyle
   \mathcal{A} \;=\;
   \sum_{i=1}^{S}
   (\mathcal{Q}_s \mathcal{X})_i
   \odot
   \operatorname{Softmax}((\mathcal{Q}_s \mathcal{X}))_i
   \in \mathbb{R}^{B \times N \times 1},
   $}
\label{equ:avg_correlation_corrected}
\end{equation}
where $\mathcal{A}$ encodes the expected spatial importance induced by the \textit{Sample Queries}.
The input features are then reweighted by $\mathcal{A}$ to obtain a query-guided representation:
\begin{equation}
\mathcal{X}_q = \mathcal{A} \odot \mathcal{X} \in \mathbb{R}^{B \times N \times C}.
\label{equ:query_aggregation}
\end{equation}
This operation highlights generalized and spatially consistent information in $\mathcal{X}$.

Accordingly, the query-guided features $\mathcal{X}_q$ are used as the input to the subspace projection
function $f_{\mathrm{sp}}(\cdot)$ defined in Eq.~\eqref{equ:subspace_projection},
yielding the feature expectation estimate required by the front-door adjustment.


\subsubsection{Feature Aggregation}
CBB outputs a combination of expected features and mediator features, capturing both generalized and task-specific information.  
First, the mediator features $\mathcal{M}$ are computed by applying a simple convolutional block to the input features $\mathcal{X}$:  
\begin{equation}
   \mathcal{M} = \mathrm{Conv}(\mathcal{X}).
\end{equation}

Next, the feature expectations 
$\hat{\mathbb{E}}[\mathcal{X}]$ (ideally, $\mathbb{E}[\mathcal{X}'] = \mathbb{E}[\mathcal{X}]$) 
and $\hat{\mathbb{E}}[\mathcal{M}]$ are estimated 
using the subspace projection $f_{\mathrm{sp}}$ described in Sec.~\ref{sec:coefficients_estimation}.  

The final output of the CBB is obtained by combining these estimates with the mediator features:
\begin{equation}
   \mathcal{F}_{\text{out}} = \hat{\mathbb{E}}[\mathcal{X}] + \hat{\mathbb{E}}[\mathcal{M}] + \mathcal{M}.
\end{equation}

Here, the first two terms implement the front-door adjustment, mitigating spurious correlations, 
while the inclusion of $\mathcal{M}$ preserves task-specific information for downstream processing. 
The CBB is fully differentiable and trained end-to-end using the downstream task loss, 
requiring no additional supervision.


